%\documentclass[12pt,oneside]{amsart}
\documentclass{scrartcl}

%\documentclass{article}

\usepackage{amsthm,amsmath,amssymb}
\usepackage[numbers]{natbib}

%\usepackage[colorlinks]{hyperref}
\usepackage{hyperref}
\hypersetup{
    colorlinks,%
    citecolor=black,%
    filecolor=black,%
    linkcolor=black,%
    urlcolor=black
}
\usepackage{hypernat}
\usepackage[top=2.5cm, bottom=2.5cm, left=2.8cm,
right=2.8cm]{geometry}
\usepackage{graphicx}


%\setlength{\parskip}{0pt}
%\setlength{\parsep}{0pt}
%\setlength{\headsep}{0pt}
%\setlength{\topskip}{0pt}
%\setlength{\topmargin}{0pt}
%\setlength{\topsep}{0pt}
%\setlength{\partopsep}{0pt}

\linespread{1}

%\usepackage{marvosym}

%\textwidth = 6.5 in \textheight = 9.2 in \oddsidemargin = 0.0 in
%\evensidemargin = 0.0 in \topmargin = -0.0 in \headheight = 0.0 in
%\headsep = 0.0in \parskip = 0.0in \parindent = 0.0in
%\def\baselinestretch{0.871}


\newtheorem{theorem}{Theorem}[section]
\newtheorem{proposition}[theorem]{Proposition}
\newtheorem{problem}[theorem]{Problem}
\newtheorem{fact}[theorem]{Fact}
\newtheorem{lemma}[theorem]{Lemma}
\newtheorem{defn}[theorem]{Definition}
\newtheorem{corollary}[theorem]{Corollary}
\newtheorem{remark}{Remark}[section]
\newtheorem{example}[theorem]{Example}

\newcommand{\E}{{\mathbb E}}
\newcommand{\se}{{\mathcal{E}}}
\newcommand{\R}{{\mathbb R}}
\renewcommand{\P}{{\mathbb P}}
\newcommand{\ac}{{\mathcal{A}}}
\newcommand{\A}{{\mathcal{A}}}
\newcommand{\B}{{\mathcal{B}}}
\newcommand{\C}{{\mathcal{C}}}
\newcommand{\M}{{\mathcal{M}}}
\newcommand{\Mf}{{\mathfrak{M}}}
\newcommand{\T}{{\mathcal{T}}}
\newcommand{\Z}{{\mathcal{Z}}}
\newcommand{\K}{{\mathcal{K}}}
\newcommand{\lc}{{\mathcal{L}}}
\newcommand{\Ps}{{\mathcal{P}}}
\newcommand{\G}{{\mathcal{G}}}
\newcommand{\pa}{{\mathring{p}}}
\newcommand{\F}{{\cal F}}
\newcommand{\samp}{{\mathcal{X}}}
\newcommand{\X}{{\mathcal{X}}}
\newcommand{\hi}{{\hat{\Phi}}}
\newcommand{\data}{{\text{data}}}

\newcounter{rcnt}[section]
\renewcommand{\thercnt}{(\roman{rcnt})}

\def\qt#1{\qquad\text{#1}}

\def\argmin{\mathop{\rm argmin}}
\def\argmax{\mathop{\rm argmax}}


\setlength{\parskip}{1.4 \medskipamount}
%\sloppy
%\linespread{1.3}

\begin{document}

\title{STAT 238 - Bayesian Statistics  \\
  Lecture Seven} 
\subtitle{Spring 2026, UC Berkeley}

\author{Aditya Guntuboyina}


\date{04 Feb 2026}

\maketitle

\section{Derivation of the Rules of Probability for Subjective
  Probability} 

In the last class, we sketched the derivation of the rules of
probability from logical consistency \textbf{without} relying on any
imagined frequency considerations for defining probability. The
argument (taken from \citet[Chapter 2]{Jaynes}) proceeded in the
following way.

We denoted the ``plausibility'' of a proposition $A$ given some
information in the form of proposition $B$ by $(A \mid B)$. We assumed
that these plausibilities take values in the set of real numbers (no
restriction to be in the interval $[0, 1]$) and that a higher value of
plausibility represents greater belief.

To deduce the usual product rule of probability, we assumed the
existence of a continuous coordinate-wise increasing function $F$ of
two real variables such that
\begin{equation*}
  (AB \mid C) = F \left((B \mid C), (A \mid BC) \right) 
\end{equation*}
for all $A, B, C$. This function can then be employed in two different
sequences to calculate $(ABC \mid D)$ for four propositions $A, B, C, D$:
\begin{align*}
&  (ABC \mid D) =   F(F((C|D), (B|CD)), (A|BCD)) \\ &
  (ABC \mid D) = F((C|D), F((B|CD), (A|BCD))). 
\end{align*}
It is therefore natural to require that the function $F$ should be such
that the right hand sides of the above two equations produce the same
answer. From this, one gets the condition:
\begin{equation*}
  F(F(x, y), z) = F(x, F(y, z)) ~~\text{for all real numbers $x, y, z$}. 
\end{equation*}
As proved in \citet[Section 2.1]{Jaynes}, the only functions $F$ which
satisfy the above equation are of the form
\begin{equation}\label{Fw}
  F(x, y) = w^{-1}(w(x) w(y))
\end{equation}
for a positive continuous increasing function $w$. Note that the
function $w$ is not unique. For example, if $w$ satisfies \eqref{Fw},
then $\tilde{w}(x) := (w(x))^{\alpha}$ (which means $\tilde{w}^{-1}(u)
= w^{-1}(u^{1/\alpha})$) also satisfies \eqref{Fw} for every $\alpha >
0$. It can be proved that there is uniqueness up to this power
transformation. 

From here, we derived the following:
\begin{enumerate}
\item $w(AB \mid C) = w(A \mid C) w(B \mid AC) = w(B \mid C) w(A \mid
  BC)$.
\item $w(A \mid C)$ always lies between 0 and 1 with
  $w(\text{impossible}) = 0$ and $w(\text{certain}) = 1$. 
\end{enumerate}
In other words, the function $w$ applied to the plausibilities leads
to quantities which satisfy the first two rules of probability.

Next the goal is to derive the sum rule. Here we first assume that
there exists a function $S: [0, 1] \rightarrow [0, 1]$  such that
\begin{equation*}
  w(A^c \mid C) = S(w(A \mid C))
\end{equation*}
for all $A$ and $C$. Note that we are working with $w(A \mid C)$
instead of the raw plausibilities $(A \mid C)$. This allows us to use
the product rule which has already been derived. We can assume that
$S$ is strictly decreasing, $S(0) = 0$, $S(1) = 1$ and that $S$ is
continuous. 

To set up the
characterizing equation for $S(\cdot)$, consider the setting of Figure
\ref{fig}. We are looking at $A$ and $B$ such that $B^c$ is contained
in $A$ (or equivalently, $A^c$ is contained in $B$). 

\begin{figure}[htbp]
\centerline{\includegraphics[scale = 0.3]{SumRule.jpeg}}
\caption{Setting for deriving the Sum Rule}
\label{fig}
\end{figure}

In this setting, there are two different ways of
calculating the plausibility of the proposition $R = A B$ in terms of
$x := w(A)$ and $y := w(B)$ and the function $S$. Both these
calculations use the product rule. The first method for calculating
$w(R) = w(AB)$ is:
\begin{align*}
  w(AB) &= w(A) w(B \mid A) \\
        &= w(A) S \left(w(B^c \mid A) \right) \\
        &= w(A) S \left(\frac{w(B^c)}{w(A)} \right) \\
                                                      &=  w(A) S
                                                        \left(\frac{S(w(B))}{w(A)}
                                                        \right) = x S
                                                        \left(\frac{S(y)}{x}
                                                        \right). 
\end{align*}
The second method for calculating $w(R) = w(A B)$ simply switches the
roles of $A$ and $B$ in the first method:
\begin{align*}
  w(AB) &= w(B)  w(A \mid B) \\
  &= w(B) S \left(w(A^c \mid B) \right) \\
        &= w(B) S \left(\frac{w(A^c)}{w(B)} \right) \\
                                                      &=  w(B) S
                                                        \left(\frac{S(w(A))}{w(B)}
                                                        \right) = y S
                                                        \left(\frac{S(x)}{y}
                                                        \right). 
\end{align*}
It is therefore natural to assume that $S$ satisfies:
\begin{equation*}
  x S
                                                        \left(\frac{S(y)}{x}
                                                        \right) =y S
                                                        \left(\frac{S(x)}{y}
                                                        \right).  
\end{equation*}
Recall that here $x = w(A)$  and $y = w(B)$. The setting is such that
$x$ and $y$ cannot be completely arbitrary. Indeed because $B^c
\subseteq A$, we must have
\begin{equation*}
  w(B^c) \leq w(A) ~~ \text{ or equivalently } ~~ S(y) \leq x.  
\end{equation*}
Our condition on $S$ is therefore
\begin{equation*}
  xS\left(\frac{S(y)}{x} \right) = y S \left(\frac{S(x)}{y} \right)
~~  \text{for all $0 \leq x, y \leq 1$ with $0 \leq S(y) \leq x$}. 
\end{equation*}
It is now proved in \citet[Section 2.2]{Jaynes} that the above condition implies that
\begin{equation*}
  S(x) = (1 - x^{\alpha})^{1/\alpha} ~~\text{for $x \in [0, 1]$}
\end{equation*}
for some $\alpha > 0$. This $\alpha$ is uniquely determined by $S$. 

We have thus proved that
\begin{equation*}
  w(A^c|C) = \left(1 - w^{\alpha}(A|C) \right)^{1/\alpha} 
\end{equation*}
which is equivalent to
\begin{equation*}
  w^{\alpha}(A^c|C) = 1 - w^{\alpha}(A|C) ~~ \text{ or } ~~
  w^{\alpha}(A^c|C) + w^{\alpha}(A|C) = 1. 
\end{equation*}
It can now be noted that the first two rules that are satisfied by
$w(A|B)$ are also satisfied by $w^{\alpha}(A|B)$. Thus
$w^{\alpha}(A|B)$ satisfies all the three rules 
\begin{enumerate}
\item $0 \leq w^{\alpha}(A|B) \leq 1$,  $w^{\alpha}(\text{impossible}) = 0$, and
  $w^{\alpha}(\text{certain}) = 1$, 
\item $w^{\alpha}(AB|C) = w^{\alpha}(A|C) w^{\alpha}(B|AC) =
  w^{\alpha}(B|C) w^{\alpha}(A|BC)$, and
\item $w^{\alpha}(A^c|C) + w^{\alpha}(A|C) = 1$. 
\end{enumerate}
Denoting $w^{\alpha}$ by $\P$, we have
\begin{enumerate}
\item $0 \leq \P(A|B) \leq 1$,  $\P(\text{impossible}) = 0$, and
  $\P(\text{certain}) = 1$, 
\item \textbf{Product Rule}: $\P(AB|C) = \P(A|C) \P(B|AC) =
\P(B|C) \P(A|BC)$, and
\item \textbf{Sum Rule}: $\P(A^c|C) + \P(A|C) = 1$. 
\end{enumerate}
We have thus derived the rules of probability without invoking any
relationship between probability and long-run frequency.   

To summarize: If we argue in terms of plausibilities but we take care
to ensure some natural constraints for logical consistency, then we
cannot manipulate the plausibilities arbitrarily but have to reason
according to the usual rules of probability after an appropriate
transformation (this transformation is given by the function
$w^{\alpha}$).

The sum rule of probability is usually stated as
\begin{equation}\label{usum}
  \P(A \cup B|C) = \P(A|C) + \P(B|C) - \P(A B|C)
\end{equation}
where $A \cup B$ denotes the proposition ``at least one of $A$ and $B$
is true'' (Jaynes uses the notation $A + B$ for $A \cup
B$). \eqref{usum} can be derived from the stated rules as 
\begin{align*}
  \P(A \cup B|C) &= 1 - \P((A\cup B)^c|C) \\
               &= 1 - \P(A^c \cap B^c|C) \\
               &= 1 - \P(B^c | A^cC) \P(A^c|C) \\
               &= 1 - \left(1 - \P(B|A^cC) \right) \P(A^c|C) \\
               &= 1 - \P(A^c|C) + \P(B|A^cC) \P(A^c|C) \\
               &= \P(A|C) + \P(A^c B|C) \\
               &= \P(A|C) + \P(B|C) \P(A^c | BC) \\
  &= \P(A|C) + \P(B|C) (1 - \P(A | BC)) \\ &= \P(A|C) + \P(B|C) -
                                             \P(B|C) \P(A|BC) =
                                             \P(A|C) + \P(B|C) - \P(A
                                             B|C). 
\end{align*}


\begin{thebibliography}{99}

  
%\bibitem[Fisher(1934)]{Fisher1934}
%R.~A. Fisher,
%\newblock ``Probability, likelihood and quantity of information in the logic of
%uncertain inference,''
%\newblock \emph{Proceedings of the Royal Society of London. Series~A,
%Containing Papers of a Mathematical and Physical Character}, vol.~146,
%no.~856, pp.~1--8, 1934.

     \bibitem[Jaynes(2003)]{Jaynes} Jaynes, E. T, (2003)
  \textit{Probability theory: the logic of science}. Cambridge
  University Press, 2003.

%  \bibitem[Sinai(1992)]{Sinai} Sinai, Y. G, (1992)
%  \textit{Probability theory: an introductory course}. Springer
%  Verlag, 1992.  

%\bibitem[{deGroot(1986)}]{DeGrootBlackwell}DeGroot, M. H. (1986)
%       A conversation with David Blackwell.
%\newblock \emph{Statistical Science}, \textbf{1}, 40--53.

%         \bibitem[Mosteller(1987)]{Mosteller} Mosteller, F, (1987)
%  \textit{Fifty challenging problems in probability with
%    solutions}. Courier Corporation, 1987.

%    \bibitem[MacKay(2003)]{MacKay} MacKay, D. J. C., (2003)
%  \textit{Information theory, inference, and learning
%  algorithms}. Cambridge University Press, 2003.

%\bibitem[{Lindley(2013)}]{Lindley}Lindley, D. V. (2013)
%   \textit{Understanding Uncertainty}. John Wiley and sons, 2013.


  \end{thebibliography}







\end{document}

%%% Local Variables:
%%% mode: latex
%%% TeX-master: t
%%% End:
